\chapter{Benchmarking of Time-Series Databases}\label{sec:benchmarking}

Every project must choose a technology stack. This choice should be driven by
the project's requirements and objective criteria; therefore, benchmarking is a
crucial part of the decision process. This chapter presents the benchmarking
procedure for time-series databases and the results obtained for the AirScout
project.

\section{Benchmarking techniques}

Benchmarks such as ClickBench, TPC-H, and TPC-DS measure analytical database
performance, but they do not accurately represent real-time
analytics~\cite{Benchmarking:03}. They often focus on batch processing and do
not capture the unique characteristics of time-series data, such as high write
throughput and efficient querying of recent data. The solution is to use
benchmarks specifically designed for time-series databases, such as TSBS (Time
Series Benchmark Suite)~\cite{TSBS:01} and InfluxDB's own benchmarking tools.
These benchmarks simulate real-world workloads, including high write rates and
complex queries on recent data~\cite{QuestDB:01}.

\section{Benchmark via TSBS for AirScout}

The AirScout project uses the TSBS benchmark to evaluate the performance of the
databases MongoDB, InfluxDB, and TimescaleDB\@. TSBS is a comprehensive
benchmarking tool that simulates real-world time-series workloads, making it an
ideal choice for our project. It allows us to generate synthetic datasets and
execute a variety of queries to measure both read and write performance across
the different databases.

To keep the comparison as fair as possible, we used the same benchmark period,
scale factor, and random seed for all systems. However, TSBS still requires
database-specific output formats and loaders, so the effective number of loaded
metrics is not guaranteed to be bit-for-bit identical across databases. The
results should therefore be interpreted as a controlled comparison under the
same workload configuration, not as a byte-identical replay into three engines.

All benchmarks run on my laptop, which has an AMD Ryzen™ AI 9 365 w/ Radeon™
880M x 20, 32 GB RAM, and a 1 TB SSD\@.

The benchmarking process involves the following steps:

\begin{itemize}
    \item \textbf{Data Generation:} TSBS pre-generates synthetic time-series datasets that mimic real-world workloads. Using a fixed random seed makes the process deterministic and reproducible across all tested databases.
    \item \textbf{Query Generation:} Queries are also pre-generated before the benchmark run to ensure that on-the-fly generation does not affect the results. The same seed is used to guarantee identical queries across databases.
    \item \textbf{Database Setup:} Each database is set up according to best practices, ensuring optimal performance for the benchmark.
    \item \textbf{Data Loading:} The pre-generated dataset is inserted into each database while TSBS records write performance metrics such as ingestion throughput and latency.
    \item \textbf{Query Execution:} The pre-generated queries are executed against the loaded data to measure read performance, including query latency and throughput.
\end{itemize}
~\cite{TSBS:01}

\subsection{Data Generation}

TSBS generates synthetic time-series data that simulates real-world workloads.
The data includes multiple metrics, tags, and timestamps, creating a realistic
environment for benchmarking. The tool offers several parameters to customize
the data generation process, such as the number of metrics, tags, and the time
range. For the AirScout project, we configured TSBS to generate a dataset that
closely resembles the expected workload of our application, ensuring that the
benchmark results are relevant and actionable.

The following command is a template showing how the dataset was generated for
benchmarking. The database placeholder must be replaced with the concrete
target format:

\begin{lstlisting}[language=bash, caption={Template command to generate synthetic time-series data using TSBS.}, label={lst:tsbs-data-generation}]
/tsbs/tsbs_generate_data \
    --use-case="iot" \
    --seed=123 \
    --scale=100 \
    --timestamp-start="2024-01-01T00:00:00Z" \
    --timestamp-end="2024-01-15T00:00:00Z" \
    --log-interval="10s" \
    --format="[timescaledb | influx | mongo]" | gzip > data/[db]-data.gz
\end{lstlisting}

Explanation of the command parameters:

\begin{itemize}
    \item \texttt{--use-case='iot'}: Specifies the use case for data generation, in this case, Internet of Things (IoT) data.
    \item \texttt{--seed=123}: Sets a fixed random seed to ensure that the generated data is the same across different runs and databases.
    \item \texttt{--scale=100}: Determines the scale of the dataset, which can be adjusted based on the expected workload.
    \item \texttt{--timestamp-start} and \texttt{--timestamp-end}: Define the time range for the generated data, simulating a realistic time-series dataset.
    \item \texttt{--log-interval='10s'}: Specifies the interval at which data points are generated, creating a consistent time-series pattern.
    \item \texttt{--format='[timescaledb | influx | mongo]'}: Indicates the output format for the generated data, allowing for compatibility with different databases.
\end{itemize}

\subsection{Query Generation}

Similar to data generation, TSBS also pre-generates queries before the
benchmark run. This approach ensures that the query generation process does not
interfere with the performance measurements. By using the same fixed random
seed for query generation, we guarantee that the same set of queries is
executed against each database, allowing for a fair comparison of read
performance.

The following command is likewise a template for query generation. The database
placeholder must be replaced with the concrete target format:

\begin{lstlisting}[language=bash, caption={Template command to generate queries using TSBS.}, label={lst:tsbs-query-generation}]
./tsbs/tsbs_generate_queries \
    --use-case="iot" \
    --seed=123 \
    --scale=100 \
    --timestamp-start="2024-01-01T00:00:00Z" \
    --timestamp-end="2024-01-15T00:00:00Z" \
    --queries=1000 \
    --query-type="airscout-crossing" \
    --format="[timescaledb | influx | mongo]" | gzip > data/queries/[db]-queries.gz
\end{lstlisting}

Explanation of the command parameters:

\begin{itemize}
    \item \texttt{--use-case='iot'}: Specifies the use case for query generation, matching the data generation use case.
    \item \texttt{--seed=123}: Sets the same fixed random seed as used in data generation to ensure that the same queries are generated for each database.
    \item \texttt{--scale=100}: Matches the scale of the dataset to ensure that the queries are relevant to the generated data.
    \item \texttt{--timestamp-start} and \texttt{--timestamp-end}: Define the time range for the queries, ensuring they target the same period as the generated data.
    \item \texttt{--queries=1000}: Specifies the number of queries to generate, which can be adjusted based on testing needs.
    \item \texttt{--query-type='airscout-crossing'}: Selects the custom query type implemented for the AirScout use case. This query type retrieves time-series measurements within the defined benchmark time window, closely matching the access patterns of the AirScout application.
    \item \texttt{--format='[timescaledb | influx | mongo]'}: Ensures that the generated \\
          queries are compatible with the respective databases being benchmarked.
\end{itemize}

\section{Database Setup}

To simplify the setup process, we used Docker to create isolated environments
for every database. This approach allows us to easily manage and configure the
databases without affecting the host system.

\section{Data Loading and Query Execution}

After setting up the databases, we loaded the pre-generated dataset into each
database while TSBS recorded write performance metrics such as ingestion
throughput and latency. Once the data was loaded, we executed the pre-generated
queries against the loaded data to measure read performance, including query
latency and throughput.

The following command is used to load the data into TimescaleDB\@:

\begin{lstlisting}[language=bash, caption={Command to load data into TimescaleDB using TSBS.}, label={lst:tsbs-data-loading}]
cat data/timescaledb-data.gz | gunzip | ./tsbs/tsbs_load_timescaledb \
    --postgres="host=localhost user=postgres password=password sslmode=disable" \
    --pass="password" --workers=4 --batch-size=10000
\end{lstlisting}

The output of the data loading process provides key performance metrics, such
as ingestion throughput and total load time.

\begin{lstlisting}[language=bash, caption={Output of the data loading process.}, label={lst:tsbs-data-loading-output}]
loaded 174247712 metrics in 26.628sec with 4 workers
(mean rate 6543739 metrics/sec)
\end{lstlisting}

Explanation of the loading command parameters:

\begin{itemize}
    \item \texttt{--postgres}: Specifies the connection string for the TimescaleDB instance, including host, user, password, and SSL mode.
    \item \texttt{--pass}: Provides the password for authentication with the database.
    \item \texttt{--workers=4}: Indicates the number of worker threads to use for loading data, which can improve performance by parallelizing the load process.
    \item \texttt{--batch-size=10000}: Sets the number of records to insert in each batch, which can help optimize performance by reducing the number of individual insert operations.
\end{itemize}

The \texttt{cat} and \texttt{gunzip} commands are used to read the compressed
data file and pipe it into the TSBS loading tool, allowing for efficient data
loading without needing to decompress the entire file beforehand.

\subsection{Data Loading Results}

The following table summarizes the data loading results for each database:

\begin{table}[ht!]
    \centering
    \begin{tabular}{|l|r|r|r|} % chktex 44
        \hline % chktex 44
        \textbf{Database} & \textbf{Metrics Loaded} & \textbf{Load Time (sec)} & \textbf{Mean Rate (metrics/sec)} \\
        \hline % chktex 44
        TimescaleDB       & 174,247,712             & 26.628                   & 6,543,739                        \\
        InfluxDB          & 184,060,164             & 33.854                   & 5,436,870                        \\
        MongoDB           & 173,268,014             & 28.812                   & 6,013,684                        \\
        \hline % chktex 44
    \end{tabular}
    \caption{Data loading performance results for each database using TSBS.}\label{tab:tsbs-data-loading-results}
\end{table}

TimescaleDB achieved the highest ingestion rate at 6,543,739 metrics/sec,
followed by MongoDB at 6,013,684 metrics/sec, while InfluxDB had the lowest
throughput at 5,436,870 metrics/sec despite loading the most metrics overall.
Because the total loaded metric counts differ slightly between the database-
specific TSBS pipelines, these values should be read as indicative benchmark
results under the same configuration rather than as a perfectly row-identical
comparison.

\begin{tcolorbox}[
        colback=red!5!white, % Background color (light red)
        colframe=red!75!black, % Frame color (dark red)
        title=\textbf{Disclaimer}, % Title in the box header
        fonttitle=\bfseries\sffamily, % Font style for the title
        boxsep=5pt, % Space between the box and surrounding text
        arc=4pt, % Rounded corners
        boxrule=0.5pt % Thickness of the frame
    ]
    The benchmark tool TSBS is currently maintained by Timescale, but it was originally created by InfluxData\@.
\end{tcolorbox}

\subsection{Query Execution Results}

The pre-generated queries are executed against each database using 4 parallel
workers, measuring read performance and access distribution across all three
systems.

\subsubsection{TimescaleDB}

\begin{lstlisting}[language=bash, caption={Command to execute queries against TimescaleDB using TSBS.}, label={lst:tsbs-query-timescaledb}, float]
cat data/queries/timescaledb-queries.gz | gunzip | ./tsbs/tsbs_run_queries_timescaledb \
    --postgres="host=localhost user=postgres password=password sslmode=disable" --workers=4
\end{lstlisting}

Result (1000 queries):

\begin{lstlisting}[language=bash, caption={Query execution output for TimescaleDB.}, label={lst:tsbs-query-timescaledb-output}, float]
Overall query rate: 14080 queries/sec
Mean latency:    0.18 ms
Max latency:    22.16 ms
\end{lstlisting}

\subsubsection{InfluxDB}

\begin{lstlisting}[language=bash, caption={Command to execute queries against InfluxDB using TSBS.}, label={lst:tsbs-query-influxdb}]
cat data/queries/influxdb-queries.gz | gunzip | ./tsbs/tsbs_run_queries_influx \
    --urls="http://localhost:8086" --workers=4
\end{lstlisting}

Result (1000 queries):

\begin{lstlisting}[language=bash, caption={Query execution output for InfluxDB.}, label={lst:tsbs-query-influxdb-output}]
Overall query rate: 6795 queries/sec
Mean latency:    0.57 ms
Max latency:     7.48 ms
\end{lstlisting}

\subsubsection{MongoDB}

\begin{lstlisting}[language=bash, caption={Command to execute queries against MongoDB using TSBS.}, label={lst:tsbs-query-mongodb}]
cat data/queries/mongodb-queries.gz | gunzip | ./tsbs/tsbs_run_queries_mongo \
    --url="mongodb://localhost:27017" --workers=4
\end{lstlisting}

Result (1000 queries):

\begin{lstlisting}[language=bash, caption={Query execution output for MongoDB.}, label={lst:tsbs-query-mongodb-output}]
Overall query rate: 3844 queries/sec
Mean latency:    1.03 ms
Max latency:     2.85 ms
\end{lstlisting}

\subsubsection{Summary}

The following table summarizes the query execution results for all three
databases:

\begin{table}[ht!]
    \centering
    \begin{tabular}{|l|r|r|r|} % chktex 44
        \hline % chktex 44
        \textbf{Database} & \textbf{Query Rate (q/s)} & \textbf{Mean Latency (ms)} & \textbf{Max Latency (ms)} \\
        \hline % chktex 44
        TimescaleDB       & \textbf{14,080}           & \textbf{0.18}              & 22.16                     \\
        InfluxDB          & 6,795                     & 0.57                       & 7.48                      \\
        MongoDB           & 3,844                     & 1.03                       & \textbf{2.85}             \\
        \hline % chktex 44
    \end{tabular}
    \caption{Query execution performance results for each database using TSBS (1000 queries, 4 workers). Best values per category are \textbf{bold}.}\label{tab:tsbs-query-results}
\end{table}

TimescaleDB delivered the highest query throughput at 14,080~queries/sec with a
mean latency of only 0.18~ms, making it the clear leader in read performance.
InfluxDB achieved 6,795~queries/sec at 0.57~ms mean latency, while MongoDB was
slowest at 3,844~queries/sec with a mean latency of 1.03~ms. Notably, MongoDB
exhibited the lowest maximum latency (2.85~ms), whereas TimescaleDB's maximum
latency of 22.16~ms suggests occasional outlier queries under parallel load.

The performance results indicate that TimescaleDB offers superior read and
write performance for time-series workloads compared to InfluxDB and MongoDB\@.
However, the choice of database should also consider other factors such as ease
of use, scalability, and specific feature requirements beyond raw performance
metrics.

\subsubsection{Decision Matrix for AirScout}

To make the final database choice transparent, we combined the benchmark
results with project-specific criteria in a weighted decision matrix. The
scores range from 1 (poor fit) to 5 (excellent fit). The weighting reflects the
needs of AirScout, where query performance is especially important, but
ingestion rate, storage efficiency, scalability, and overall cost/flexibility
must also be considered. For TimescaleDB, InfluxDB, and MongoDB, the benchmark
results directly inform the performance scores. Prometheus and OpenTSDB are not
included in the benchmark and also not evaluated in the decision matrix because
they are more specialized for monitoring workloads and do not align well with
AirScout's broader requirements of a general-purpose time-series database that
also supports data models like relational or document-based structures.

\begin{table}[ht!]
    \centering
    \scriptsize
    \resizebox{\textwidth}{!}{
        \begin{tabular}{|l|c|c|c|c|} % chktex 44
            \hline % chktex 44
            \textbf{Criterion}      & \textbf{Weight} & \textbf{InfluxDB} & \textbf{TimescaleDB} & \textbf{MongoDB} \\
            \hline % chktex 44
            Ingestion performance   & 20\%            & 4                 & 5                    & 4                \\
            \hline % chktex 44
            Query performance       & 35\%            & 4                 & 5                    & 2                \\
            \hline % chktex 44
            Storage efficiency      & 15\%            & 5                 & 4                    & 3                \\
            \hline % chktex 44
            Scalability             & 15\%            & 4                 & 4                    & 4                \\
            \hline % chktex 44
            Costs \& flexibility    & 15\%            & 3                 & 5                    & 3                \\
            \hline % chktex 44
            \textbf{Weighted total} & \textbf{100\%}  & \textbf{4.00}     & \textbf{4.70}        & \textbf{3.00}    \\
            \hline % chktex 44
        \end{tabular}}
    \caption{Weighted decision matrix for selecting a time-series database for AirScout. Scores are project-specific and use a scale from 1 (poor fit) to 5 (excellent fit).}\label{tab:airscout-decision-matrix}
\end{table}

The matrix confirms the interpretation of the raw benchmark results. InfluxDB\@
scores strongly in storage efficiency and remains a capable specialized TSDB\@.
MongoDB\@ benefits from schema flexibility but trails the specialized
time-series systems in query performance. TimescaleDB\@ achieves the highest
overall score because it offers the strongest measured read and write
performance while also providing SQL, PostgreSQL integration, and a flexible
path for future analytics.

In the end, we chose TimescaleDB for AirScout due to its excellent performance
in both data loading and query execution, as well as its seamless integration
with PostgreSQL, the database we already knew and trusted, which allows us to
leverage existing tools and expertise while benefiting from time-series
optimizations.
